\section*{Exercises on \S\arabic{exercisegroup}}
\phantomsection
\addcontentsline{toc}{section}{Exercises on \protect\S\arabic{exercisegroup}}

\begin{enumerate}
\def\labelenumi{\arabic{enumi})}
\tightlist
\item
  Show that for \(x\) real tending to \(+\infty\) the function
\end{enumerate}

\[
f (x) = \left(x \cos^ {2} x + \sin^ {2} x\right) e ^ {x ^ {2}}
\]

is monotone, but is neither comparable to \(e^{x^2}\) nor weakly comparable to \(x^{1/2}e^{x^2}\) .

\begin{enumerate}
\def\labelenumi{\arabic{enumi})}
\setcounter{enumi}{1}
\tightlist
\item
  Let \(\varphi\) be a strictly positive function, defined and increasing for \(x > 0\) , and such that \(\varphi \gg 1\) .
\end{enumerate}

\begin{enumerate}
\def\labelenumi{\alph{enumi})}
\item
  Show that if the function \(\log \varphi(x) / \log x\) is increasing, then the relation \(f \ll g\) between functions \(>0\) implies that \(\varphi \circ f \ll \varphi \circ g\) if \(g \gg 1\) .
\item
  Give an example where \(\log \varphi(x)/\log x\) is decreasing, f and g are two functions \textgreater0 such that \(g \gg 1\) and \(f \sim g\) , but \(\varphi \circ f\) is not equivalent to \(\varphi \circ g\) .
\end{enumerate}

3) a) Let \((\alpha_{i},\beta_{i})\) be \(n\) distinct pairs of real numbers \(\geqslant 0\) , distinct from \((0,0)\) , and such that \(\inf_{i}\alpha_{i} = \inf_{i}\beta_{i} = 0\) . Consider the equation

\[
f (x, y) = \sum_ {i = 1} ^ {n} a _ {i} x ^ {\alpha_ {i}} y ^ {\beta_ {i}} \left(1 + \varphi_ {i} (x, y)\right) = 0
\]

where the \(a_{i}\) are real numbers \(\neq0\) and the \(\varphi_{i}\) are continuous functions on a square \(0\leqslant x\leqslant a,0\leqslant y\leqslant a\) , tending to 0 as \((x,y)\) tends to \((0,0)\) . Suppose that there exists a function g which is positive and continuous on an interval {[}0,b{]}, tending to 0 with x and such that \(f(x,g(x))=0\) for all \(x\in[0,b]\) . For every real number \(\mu>0\) show that \(g(x)/x^{\mu}\) tends to a limit, finite or infinite, as x tends to 0 (use V, p.~217, prop. 9, considering, for every number \(t\geqslant0\) , the equation \(f(x,tx^{\mu})=0\) ).

\begin{enumerate}
\def\labelenumi{\alph{enumi})}
\setcounter{enumi}{1}
\item
  For \(g(x)/x^{\mu}\) to tend to a finite limit \(\neq0\) it is necessary that \(\mu\) should be such that there exist at least two distinct pairs \((\alpha_{h},\beta_{h})\) and \((\alpha_{k},\beta_{k})\) such that \(\alpha_{h}+\mu\beta_{h}=\alpha_{k}+\mu\beta_{k}\) and such that, for every other pair \((\alpha_{i},\beta_{i})\) , one has \(\alpha_{i}+\mu\beta_{i}\geqslant\alpha_{h}+\mu\beta_{h}\) . One thus obtains a finite number of possible values \(\mu_{j}(1\leqslant j\leqslant r)\) ; the numbers \(-1/\mu_{j}\) are the gradients of the affine lines in the plane \(R^{2}\) which contain at least two of the set of points \((\alpha_{i},\beta_{i})\) and such that all the other points of this set lie above the line considered (``Newton polygon'').
\item
  Let \(\mu_{1}\) be the smallest of the numbers \(\mu_{j}\) . Show that \(g(x) / x^{\mu_1}\) tends to a finite limit (possibly zero). (Show first that one may always assume that if \(i\) and \(j\) are two distinct indices one cannot have both \(\alpha_{i} \leqslant \alpha_{j}\) and \(\beta_{i} \leqslant \beta_{j}\) ; deduce that one can assume \(\alpha_{1} = 0\) , \(\alpha_{i} > 0\) and \(\beta_{i} < \beta_{1}\) for \(i \neq 1\) ; now, putting \(g(x) = t(x)x^{\mu_1}\) , show that \(t(x)\) cannot tend to \(+\infty\) as \(x\) tends to 0.)
\end{enumerate}

\(d)\) Deduce from \(c\) ), by recursion on \(r\) , that there exists an index \(j\) such that \(g(x) / x^{\mu_j}\) tends to a finite limit \(\neq 0\) as \(x\) tends to 0.

\setcounter{exercisegroup}{2}
\refstepcounter{exercisegroup}
